% 本文件由 LaTeX 排版 Agent 根据有效 translation.json 自动生成。
% author=Apocrypha; work_id=0819; title=安德肋行传

\chapter{安德肋行传}

% structure: p0001 source=h1 target=omitted reason=page-title-already-rendered-by-parent

我们众人，即\Place{阿哈雅}众教会的司铎和执事，亲眼所见，如今写信给所有因\Person{耶稣基督}之名建立的教会，无论在东方西方、南方北方。愿平安归于你们，以及所有相信唯一天主、完美天主圣三、真实无生之父、真实独生子、真实发自父并寓于子的圣神的人，为彰显一位在父与子内具有珍贵神性的圣神。这信德我们是从我们的主耶稣基督的宗徒、圣\Person{安德肋}所学到的，我们亲眼目睹了他的蒙难，便毫不迟疑地按照我们的能力将其叙述出来。\par

因此，总督\Person{埃该阿特斯}来到\Place{帕特雷}城，开始强迫信奉基督的人崇拜偶像。圣\Person{安德肋}跑上前去，对他说：\Quote{你身为审判世人的人，理当认识你那天上的审判者，认识祂后就敬拜祂；敬拜那真天主，就应转离那些不是真神的偶像。}\par

\Person{埃该阿特斯}对他说：\Quote{你就是那拆毁众神庙宇、劝人信奉那新近出现、罗马皇帝下令禁止的宗教的安德肋吗？}\par

圣\Person{安德肋}说：\Quote{罗马皇帝从未认识真理。而天主子为拯救人类而来，清楚地教导：这些偶像不仅不是神，更是最可耻的魔鬼，与人类为敌，教导人得罪天主，以致天主因被得罪而转身不听；因祂转身不听，他们就被魔鬼束缚；魔鬼甚至驱使他们到这般地步，以致当他们离开身体时，被发现孤零零、赤身露体，除了罪以外，什么也没有带走。}\par

\Person{埃该阿特斯}说：\Quote{这些都是多余而虚妄的话。至于你的耶稣，因向犹太人宣讲这些事，他们便把他钉在十字架上。}\par

圣\Person{安德肋}回答说：\Quote{唉！你若能认识十字架的奥秘，那人类生命的创造者为恢复我们而以何等理性的爱，甘愿忍受这十字架，并非不愿，而是甘心！}\par

\Person{埃该阿特斯}说：\Quote{既然他被自己的门徒出卖，被犹太人抓住，带到总督面前，又按他们的要求被总督的士兵钉在架上，你怎么说他甘愿忍受十字架呢？}\par

圣\Person{安德肋}说：\Quote{正因如此，我说他甘愿，因为当祂被门徒出卖时，我正与祂同在。因为在祂被出卖之前，祂对我们说过，祂将被出卖并被钉在十字架上，为拯救人类，并预言祂第三日要复活。我的兄弟\Person{伯多禄}对祂说：\BibleRef{玛 16:22}\BibleQuote{主，千万不可！这事绝不会临到你身上。}祂就生气，对\Person{伯多禄}说：\InnerQuote{撒殚，退到我后面去！因为你不思念天主的事。}为了更充分地说明祂是甘愿受难的，祂对我们说：\BibleRef{若 10:18}\BibleQuote{我有权舍弃我的生命，也有权再取回它。}最后，当祂与我们一同用晚餐时，祂说：\BibleRef{玛 26:21}\BibleQuote{你们中间有一个人要出卖我。}因此，这些话使众人都极其忧伤，为使猜疑无可置疑，祂指明说：\InnerQuote{我递饼给谁，谁就是出卖我的人。}于是，当祂把饼递给一位同门时，就将未来的事当作已发生的事来说明，显示出祂是心甘情愿被出卖的。因为祂没有逃跑，使出卖者落空；反而留在那祂知道出卖者会来的地方，等候着他。}\par

\Person{埃该阿特斯}说：\Quote{我奇怪，你这样一个有理智的人，竟无论如何也要为他辩护；因为无论是甘愿还是不愿意，你都得承认他被钉在十字架上。}\par

圣\Person{安德肋}说：\Quote{这就是我所说的，如果你现在能理解的话：十字架的奥秘是伟大的，你若愿意听，就如所宜，请听我说。}\par

\Person{埃该阿特斯}说：\Quote{这不能称为奥秘，而是刑罚。}\par

圣\Person{安德肋}说：\Quote{这刑罚正是人类恢复的奥秘。你若稍加留意听，就会明白。}\par

\Person{埃该阿特斯}说：\Quote{我固然会耐心听；但除非你顺从地服从我，否则你将在自己身上领受十字架的奥秘。}\par

圣\Person{安德肋}回答说：\Quote{我若害怕十字架，就不会宣讲十字架的荣耀了。}\par

\Person{埃该阿特斯}说：\Quote{你的话是愚昧的，因为你宣称十字架不是刑罚，而且你凭着鲁莽不怕死亡的刑罚。}\par

圣\Person{安德肋}说：\Quote{我并非凭鲁莽，而是凭信德，不怕死亡的刑罚；因为罪的死亡是严苛的。正因如此，我希望你听十字架的奥秘，好叫你也许能认识而相信，相信而终究达到你灵魂的更新。}\par

\Person{埃该阿特斯}说：\Quote{显示为已毁坏的东西是要更新的。你的意思是我的灵魂已毁坏，好叫我能通过你所讲的那种信德达到灵魂的更新吗？}\par

真福安德肋回答说：\Quote{这正是我渴望时间来学习的，也是我将要教导和显明的：虽然人的灵魂被摧毁，但通过十字架的奥迹，它们将得到更新。因为第一个人因悖逆之树带来了死亡；而人类有必要通过那树的苦难，使进入世界的死亡被驱逐出去。既然第一个人——他通过那树的悖逆将死亡带入世界——是从无瑕的土地中产生的，那么天主子就必须从无瑕的童贞女中生成为一个完美的人，为要恢复人类因亚当而失去的永生，并藉着十字架之树砍断肉欲之树。祂悬挂在十字架上，为那毫无节制地伸出的双手伸出了自己无瑕的双手；为那禁树最甜美的食物，祂接受了苦胆作为食物；祂将我们的必死性披在自己身上，并将祂的不朽性赐予我们。}\par

厄该雅特说：\Quote{用这些话你只能引导那些相信你的人；但除非你来是答应我献祭给全能的神祇，否则我将命令你，在鞭打之后，被钉在你所赞扬的那十字架上。}\par

真福安德肋说：\Quote{我每天向唯一真天主、全能的天主祭献；不是馨香的烟，也不是吼叫的公牛的肉，也不是山羊的血，而是每天在十字架的祭台上祭献一只无瑕的羔羊；虽然所有信众都分享祂的身体并饮祂的血，但这被祭献的羔羊此后仍然是完整而活的。因此，祂确实被祭献，祂的身体也确实被信众所吃，祂的血也同样被喝；然而，正如我所说，祂仍然是完整的、无瑕的、活的。}\par

厄该雅特说：\Quote{这怎么可能？}\par

真福安德肋说：\Quote{你若想知道，就取门徒的样式，好让你能学到你所询问的事。}\par

厄该雅特说：\Quote{我将通过酷刑从你那里逼出这知识的恩赐。}\par

真福安德肋宣称：\Quote{我感到惊奇，你这样一个聪明的人，竟会陷入这样的愚昧，以为通过你的酷刑就能说服我向你揭示天主的神圣事物。你已经听到了十字架的奥迹，你已经听到了祭献的奥迹。如果你相信被钉十字架的基督、天主子，那么我将完全向你揭示，那被宰杀的羔羊，在被祭献并吃了之后，如何能活着，在祂的王国中保持完整无瑕。}\par

厄该雅特说：\Quote{那么，那羔羊被宰杀并被所有百姓吃了之后，如何能留在他的王国中，如你所说？}\par

真福安德肋说：\Quote{你若全心相信，就能学到；但若不相信，你绝不能达到如此真理的观念。}\par

于是厄该雅特暴怒，下令将他关进监狱；他被关进监狱时，几乎全省的民众都聚集到他那里，想要杀死厄该雅特，并打破监狱的门释放真福安德肋宗徒。\par

真福安德肋用这些话劝诫他们，说：\Quote{不要以煽动和魔鬼般的喧嚣扰乱我们主耶稣基督的和平。因为我的主，当祂被出卖时，以极大的忍耐承受了；\InnerQuote{祂没有争辩，没有喊叫，街上也没有人听见祂喊叫。}\BibleRef{玛 12:19}因此，你们也当保持沉默、安静与和平；不要阻止我的殉道，反而要准备好自己，作为战士到主面前，以便你们能以无所畏惧的灵魂克服威吓，并以身体的坚忍胜过伤害。因为这暂时的跌倒并不可怕；那没有结局的才可怕。对人的恐惧如同烟云，上升聚集时便消失了。而那些永无终局的酷刑才当畏惧。因为那些偶然较轻的酷刑，任何人都能忍受；但若沉重，它们很快就夺去生命。但那些永久的酷刑，那里有日日夜夜的哭泣、哀恸、叹息和无尽的折磨，总督厄该雅特并不惧怕走向那里。因此，你们更当预备自己，\InnerQuote{藉暂时的苦难达到永久的安息，永远昌盛，并与基督一同为王。}\BibleRef{格后 4:17}}\par

圣宗徒安德肋用这些和类似的话劝诫了民众整夜；当白昼的光明破晓时，厄该雅特派人将他带来，命令将真福安德肋带到面前；他坐在审判台上，说：\Quote{我以为你经过一夜的反思，已经将你的心思从愚昧中转回，放弃了对基督的赞美，以便你能与我们在一起，不抛弃生活的快乐；因为为了某种目的而走向十字架的苦难，并把自己交给最可耻的刑罚和火刑，是愚昧的。}\par

圣安德肋回答说：\Quote{如果你相信基督，并抛弃偶像崇拜，我将能与你一同喜乐；因为基督派我来到这个省份，在这里我为基督获得了为数不小的子民。}\par

厄该雅特说：\Quote{正因如此，我强迫你奠酒祭神，以便这些被你欺骗的百姓能放弃你教导的虚妄，他们自己向众神献上感恩的奠祭；因为在阿哈雅，没有一个城市中他们的庙宇没有被抛弃和荒废的。现在，藉着你，让它们重新恢复对神像的敬拜，以便那些向你发怒的众神，因这件事而喜悦，可以带你回到他们和我们的友谊中。但若不如此，你将等待各种酷刑，作为众神的报复；在这之后，被钉在你所赞扬的十字架上，你将死去。}\par

圣安德肋说：\Quote{死亡之子，为永远焚烧所预备的糠秕\BibleRef{玛 3:12}啊，请听我，天主的仆人和耶稣基督的宗徒说话。至今我仁慈地与你谈论信德的圆满，好使你接受真理的阐释，成为真理的维护者而得以成全，从而藐视虚妄的偶像，敬拜那在天上的天主；但既然你最终仍留在同样的无耻中，并因你的恐吓而以为我胆怯，那么，就把你认为更厉害的酷刑加在我身上吧。因为我越是为宣认祂的名而在酷刑中忍受，就越能讨我的君王喜悦。}\par

于是总督埃格阿特斯恼怒，下令用酷刑折磨基督的宗徒。因此他被七组、每组三个士兵拉开来猛烈鞭打，然后被抬起带到不敬神的埃格阿特斯面前。埃格阿特斯对他说：\Quote{你听我说，安德肋，从你的流血中收回你的心思吧；但若你不听从我，我就要让你死在十字架上。}\par

圣安德肋说：\Quote{我是基督十字架的奴隶，我更应当祈求达到十字架的凯旋，而不是害怕；但为你却存留了永久的折磨，然而，你若相信我的基督，在考验了我的坚忍之后，你仍可逃脱它。因为我为你的丧亡而忧虑，并不为自己的受苦而烦恼。我的受苦不过占据一两天的时间；但你的折磨却永无终期。所以，从现在起，不要再增添你的痛苦，也不要点燃那永远之火。}\par

于是埃格阿特斯恼羞成怒，下令将圣安德肋钉在十字架上。安德肋离开众人，走向十字架，以清亮的声音对它说：\Quote{欢乐吧，十字架！你曾被基督的身体祝圣，并被祂的肢体如珍珠般装饰。的确，在我的主登上你之前，你曾有属地的恐惧；但如今你被赋予属天的渴慕，依照我的祈祷被预备好了。因为我知道，从那些相信的人那里，你在祂内拥有多少恩宠，预先准备了多少恩赐。于是，我无忧无虑、满心喜乐地来到你面前，好使你也欢欣地接纳我，那位曾被悬在你身上的祂的门徒；因为你对我始终忠诚，我一直渴望拥抱你。啊，良善的十字架，你从主的肢体获得了优雅和美丽；啊，备受渴望、热切向往、炽烈寻求，并已为我的灵魂预先准备好的十字架，求你把我从人间带走，将我归还给我的主，好使祂藉着你接纳我——因为祂曾藉着你救赎了我。}\par

说了这些话之后，圣安德肋站在地上，定睛望着十字架，脱下衣服交给行刑者，并劝勉弟兄们让行刑者前来执行所吩咐的事情；因为他们正站在远处。他们上前来，把他抬上十字架；用绳索将他的身体横着拉紧，只绑住他的脚，并没有伤及骨节，这是奉了总督的命令：因为他希望他在悬挂时受苦，并在夜间被悬吊着活活被狗吃掉。\par

有一大群弟兄站立在旁，将近两万人；他们看见行刑者站在远处，并没有对被悬吊者通常所受的苦施加于圣安德肋身上，以为他们会再听见他说些什么；因为他一边悬挂着，一边微笑着摇头。斯特拉托克勒问他：\Quote{安德肋，天主的仆人，你为什么微笑？你的笑使我们哀恸哭泣，因为我们失去了你。}圣安德肋回答说：\Quote{我的儿子斯特拉托克勒，我岂能不嘲笑埃格阿特斯那空洞的诡计？他以为能借此报复我们。我们与他及其计划毫无关系。他听不见；因为如果他能够听见，他必会从经验得知，属于耶稣的人是不会受罚的。}\par

说了这些话之后，他又向全体会众讲话，因为民众因埃格阿特斯不义的判决而激愤地聚集起来：\Quote{站在我身边的男男女女、孩童、长者、为奴的、自主的，以及所有愿意听的人，我恳求你们，你们因我的缘故而聚集在这里，要离弃这整个现世生命，赶快接受那引向属天事物的我的生命，彻底藐视一切暂时之物，坚定那些相信基督之人的心意。}他劝勉众人，教导说这短暂生命的苦难，不配与未来永生的赏报相比。\par

群众听到他所说的话，并没有离开那地方，圣安德肋反而更加继续向他们说更多的话。他讲了那么多，以至于持续了三天三夜，没有一个人疲倦而离开他。到了第四天，当他们看到他的高贵、才智的不倦、言语的丰富、劝勉的助益、灵魂的坚定、精神的清醒、心志的稳固和理性的圆满时，他们便对埃格阿特斯愤怒不已；所有人一致冲向审判台，向坐着的埃格阿特斯喊道：\Quote{总督，你作出什么判决？你判得不公；你的裁决是不敬神的。这人做了什么错事？他行了什么恶？全城都被搅乱了；你使我们所有人忧伤；不要出卖凯撒的城。请慷慨地赐给阿哈雅一个义人；请慷慨地赐给我们一个敬畏天主的人；不要处死一个虔诚的人。他已被悬挂四天，却仍然活着；没有吃任何东西，却使我们众人饱足。把这人从十字架上放下来，我们众人便会寻求智慧；释放这人，整个阿哈雅便会蒙受怜悯。他不该受此酷刑，因为尽管被悬挂，他仍不停宣讲真理。}\par

当总督拒绝听他们时，起初他用手势示意人群离开，但他们开始壮起胆来反对他，人数约有二万。总督见他们似乎有些疯狂，害怕可怕的事会降临到他身上，便从审判台起身，与他们一同离去，答应释放有福的\Person{安德肋}。有人先走一步，去告诉宗徒他们来到这地方的原因。\par

因此，当全体民众因有福的安德肋即将被释放而欢呼时，总督上前来，所有的弟兄也与\Person{玛克西米拉}一同喜乐。有福的安德肋听了这事，对站在旁边的弟兄们说：\Quote{当我即将归主时，对他说什么是必要的，我也要说。厄革亚特，你为什么再次来到我们这里？你与我们不相干，为什么到我们这里来？你又要胆敢做什么？又要图谋什么？告诉我们。你是改变了主意来释放我们吗？我不相信你真的改变了主意。我也不相信你说你是我的朋友。总督啊，你释放那被捆绑的人吗？绝不。因为我有一位与他永远同在；我有一位与他同活无数世代。我要去他那里；我向他奔去，他也曾使我认识你，并对我说：\InnerQuote{不要让那可怕的人恐吓你；不要以为他会抓住你，你是属于我的，因为他是你的敌人。} 因此，借着那转向我的一位，我认识了你，便从你那里被释放了。但如果你愿意相信基督，就如我向你应许的，有一条进入之路会为时间而打开；但如果你来只是为了释放我，此后我将不能活着从这十字架上被放下来。因为我和我的亲属要往我们那里去，让你仍是你的样子，而你不认识你自己。因为我已经看见我的君王，我已经朝拜他，我已经站在他面前，在那里有天使的共融，他是唯一的元首，那里有光明没有黑夜，那里花朵永不凋谢，那里没有烦恼，也听不到忧伤的名字，那里有无尽的欢乐与喜乐。噢，有福的十字架！没有对你的渴望，没有人能进入那地方。厄革亚特，我却为你的悲惨而忧虑，因为永久的灭亡正等着你。那么，为了你自己，跑吧，可怜的人，趁你还能的时候，免得你到时想要却无能为力。}\par

因此，当他试图走近十字架的木架，想要释放有福的安德肋时，全城的人都为他喝彩。圣安德肋大声说：\Quote{主，不要容许安德肋被捆绑在祢的木架上而被释放；不要把我这处于祢奥秘中的人交给无耻的魔鬼。耶稣基督，不要让祢的仇敌释放我，我是因祢的恩宠而被悬挂的；父啊，不要让这个微不足道的人再贬低那认识祢伟大的人。} 于是，刽子手们伸出手，却完全不能碰他。其他人也一次又一次试图释放他，却没有人能靠近他，因为他们的手臂麻木了。\par

然后，有福的安德肋恳请民众，说：\Quote{我恳切地请求你们，弟兄们，让我先向我的主献上一次祈祷。所以你们就着手释放我吧。} 全体民众因此因这恳请而安静下来。然后，有福的安德肋大声呼喊说：\Quote{主，不要容许祢的仆人此时从祢那里被移开；因为现在是我的身体被交于土地的时候，祢将命令我来到祢面前。祢是赐予永生者，我所爱的老师，我在这个十字架上所宣认的、我所认识的、我所拥有的祢，请收纳我，主；我既宣认了祢并听从了祢，现今在这话语上求祢垂听我；在我的身体从十字架上下来之前，请祢接纳我到祢那里，好使我的许多同族能因我的离去而得以接近祢，并在祢的尊威中找到安息。}\par

当他说完这话，在众人眼前显得喜乐欢腾；因为从天上发出如同闪电的极大光辉照耀着他，并环绕他，以至于因这光辉，肉眼完全不能看见他。这耀眼的光辉持续了约半小时。当他如此说了并更加光荣了主之后，光辉退去，他交付了灵魂，并随同光辉一同归向了主，感谢他。\par

至福的宗徒安德肋去世后，\Person{玛克西米拉}——她是显贵妇女中最有权势的，并留在那些来者中——一得知宗徒已归主，便上前与\Person{斯特拉托克勒}一起关注十字架，毫不理会在场的人，恭敬地将至福宗徒的身体从十字架上取下。到了晚上，她尽了必要的照料，用贵重的香料预备了身体并安葬，将它放在自己的坟墓里。因着\Person{厄革亚特}的残暴性情和无法无天的行为，她已与他分离，为自己选择了圣洁而安静的生活；她与基督的爱结合，同弟兄们一起有福地度过了余生。\par

\Person{厄革亚特}曾对她非常纠缠，应许让她成为他财富的主妇；但未能说服她，便极其愤怒，决定公开控告全体民众，并要送信给\Person{凯撒}，控告\Person{玛克西米拉}和全体民众。当他在众官员面前安排这些事时，深夜，他起身，不被所有人看见，被魔鬼折磨，从高处坠落，滚到城中的市场中央，断了气。\par

这事报告给他的兄弟\Person{斯特拉托克勒}；他派了仆人去，吩咐他们把这人埋在同那些横死的人一起。但他不曾寻求任何自己的财物，却说：\Quote{愿我所信的主耶稣基督，绝不许我碰触我兄弟的任何财物，免得那敢于杀害主之宗徒的人的定罪也玷污了我。}\par

这些事发生在\Place{阿哈雅}省，\Place{帕特雷}城，在\OriginalTerm{十二月朔日前一日}{kalends}，在那里他的善行至今仍被纪念，为光荣和赞美我们的主耶稣基督，愿光荣归于他，直到永远。阿们。\par

\SourceRecord{Translated by Alexander Walker.；Ante-Nicene Fathers；Vol. 8.；Edited by Alexander Roberts, James Donaldson, and A. Cleveland Coxe.；Buffalo, NY: Christian Literature Publishing Co.,；1886.；<http://www.newadvent.org/fathers/0819.htm>.}
